\section{Conclusion}

SERVE introduces a margin-aware, tree-based gate for semantic LLM caches
that achieves higher hit rates at controllable false-serve budgets. Where
every deployed semantic cache discards the runner-up similarity after
nearest-neighbor search, SERVE repurposes it, together with the margin
between the top two matches, as input to a gradient-boosted classifier
that decides whether a cached answer can be safely reused.
% E15

The margin alone accounts for 83\,\% of the total hit-rate improvement
over a similarity-only gate at a 1\,\% false-serve budget.
% E4
The gate adds 0.8\,$\mu$s of inference latency per query with 0.0\,\textmu s standard deviation across runs, and none of its four features requires
an additional embedding call or index lookup: every figure is a byproduct
of the top-$k$ search the cache already performs, making SERVE a pluggable
component compatible with any similarity-based reuse rule.
% E5
% E14
On QQP with MiniLM embeddings, SERVE improves hit rate by 11.9\,\% over a
fixed threshold and 6.5\,\% over vCache at $\beta \le 1\,\%$, and the
gain grows monotonically with query recurrence, exceeding $3\times$ the
served hits in the strictest-budget, highest-recurrence cells.
% E2
% E7

Reducing the number of LLM calls needed to satisfy a stream of queries
lowers inference cost and latency. A gate that ships as a standalone,
sub-microsecond component lowers the barrier to adopting semantic caching
in production, which in turn makes large-model deployment feasible under
tighter resource constraints.

Several directions merit further investigation. Extending the gate to
multi-turn conversations, where the reuse decision depends on dialogue
context, would test whether the margin signal retains predictive power
when the cache contains utterances rather than standalone questions.
Dynamic cache eviction policies that account for the margin distribution
of stored entries could tighten the hit-rate--reliability frontier.
Online learning of the gate, retraining incrementally as cache content
drifts, would eliminate the current one-time training requirement.

Cosine similarity alone is not a reliable signal for semantic cache
reuse. The margin between the top match and its runner-up, discarded by
every incumbent system, is the missing piece that makes the reuse
decision tractable. SERVE demonstrates that recovering this signal costs
nothing and delivers substantial, measurable improvement at the operating
points that matter for real deployments.
